\documentclass[11pt]{article}
\usepackage{manual_docs_style}
\externaldocument{build/run_graph_schema}[run_graph_schema.pdf]
\externaldocument{build/generation_policy_schema}[generation_policy_schema.pdf]
\externaldocument{build/simulation_stage}[simulation_stage.pdf]

\begin{document}

\docTitle{Legacy Model Run Specs Compatibility}
\phantomsection\label{docstart:model_run_specs_schema}

\section*{\texorpdfstring{\titlecode{model_run_specs.json}}{model_run_specs.json} Compatibility Schema}
\label{sec:model-run-specs-schema}

\code{model_run_specs.json} is a legacy compatibility artifact.
It is no longer the primary source of truth for planned runs.

The new source of truth is the resolved run graph produced by the candidate-generation stage:

\begin{DocCode}
models/simulink/<MODEL>/plans/<policy_id>/run_nodes.jsonl
models/simulink/<MODEL>/plans/<policy_id>/run_edges.jsonl
\end{DocCode}

New code should consume the resolved run graph directly.
\code{model_run_specs.json} may still be emitted during migration for older tools that expect the previous nested baseline/child/time-zero structure.

\section*{Status}

The old contract treated \code{model_run_specs.json} as the human-authored planned topology consumed by Stage: Simulate.
The new contract separates those responsibilities:

\begin{center}
\begin{tabular}{p{0.30\linewidth}p{0.60\linewidth}}
\hline
\textbf{Artifact} & \textbf{Responsibility} \\
\hline
\code{generation_policies/<policy_id>.json} & Human-authored or agent-authored sampling recipe for one candidate-generation campaign. \\
\code{plans/<policy_id>/run_nodes.jsonl} & Resolved runnable run recipes. This is the canonical source of truth for \code{run_id}, \code{recipe_hash}, and simulator inputs. \\
\code{plans/<policy_id>/run_edges.jsonl} & Resolved graph relationships among baseline, intervention, and time-zero/static-change runs. \\
\code{model_run_specs.json} & Optional legacy compatibility export. It should not be edited manually and should not be consumed by new code. \\
\hline
\end{tabular}
\end{center}

\section*{When This File Exists}

A compatibility file may be written in either of these locations:

\begin{DocCode}
models/simulink/<MODEL>/plans/<policy_id>/model_run_specs.compat.json
\end{DocCode}

or, when an old script requires the historical root-level path:

\begin{DocCode}
models/simulink/<MODEL>/model_run_specs.json
\end{DocCode}

The preferred location is \code{plans/<policy_id>/model_run_specs.compat.json}, because it makes clear which policy produced the compatibility artifact.
If the root-level \code{model_run_specs.json} exists, it must be treated as a generated compatibility copy, not as an editable source file.

\section*{Legacy Shape}

The compatibility file preserves the old nested shape:

\begin{DocCode}
{
  "<baseline_uuid>": {
    "baseline_parameters": {
      "simulation_specific_parameter_0": 1.0,
      "simulation_specific_parameter_1": 2.0
    },
    "baseline_parameters_hash": "<32-hex-hash>",
    "children": {
      "<intervention_uuid>": {
        "intervention_time": 3.0,
        "parameters": {
          "simulation_specific_parameter_0": 1.5
        },
        "parameter_hash": "<32-hex-hash>",
        "time0_baseline_hash": "<32-hex-hash>",
        "time0_baseline_uuid": "<32-hex-id>",
        "direction": "increase"
      }
    }
  }
}
\end{DocCode}

This shape is retained only so that legacy consumers can recover the same baseline/intervention/time-zero topology they previously expected.

\section*{Mapping from the Resolved Run Graph}

The compatibility exporter builds this file from \code{run_nodes.jsonl} and \code{run_edges.jsonl}.
For each \code{family_id}, the exporter should:

\begin{enumerate}
\item find the baseline node for the family;
\item create one top-level baseline entry;
\item find all intervention nodes connected to that baseline through \code{baseline_to_intervention} edges;
\item for each intervention node, find its corresponding time-zero/static-change node through an \code{intervention_to_time0_baseline} edge;
\item write the old nested child entry.
\end{enumerate}

The compatibility file must not introduce runnable runs that are absent from the resolved run graph.

\section*{Identifiers}

The new canonical identifier is \code{run_id} from \code{run_nodes.jsonl}.
The old compatibility schema uses \code{baseline_uuid}, \code{intervention_uuid}, and \code{time0_baseline_uuid}.

During compatibility export:

\begin{itemize}
\item \code{baseline_uuid} should correspond to the baseline node's \code{run_id};
\item \code{intervention_uuid} should correspond to the intervention node's \code{run_id};
\item \code{time0_baseline_uuid} should correspond to the time-zero/static-change node's \code{run_id}.
\end{itemize}

If a legacy consumer requires 32-character lowercase hexadecimal identifiers, the exporter may provide compatible aliases.
New code must not depend on these aliases; it should use \code{run_id} from the resolved run graph.

\section*{Hashes}

The new canonical identity field is \code{recipe_hash} in \code{run_nodes.jsonl}.
The old compatibility fields are \code{baseline_parameters_hash}, \code{parameter_hash}, and \code{time0_baseline_hash}.

These fields exist only for legacy consumers.
They should be generated from the canonical run recipes and should not be edited manually.
New code should compare \code{recipe_hash}, not these legacy hash fields.

\section*{Skipped Children}

The resolved run graph contains only runnable run nodes.
It does not contain fake or null-valued children.
If a requested parameter-direction pair cannot be sampled, the missing direction is recorded in:

\begin{DocCode}
plans/<policy_id>/generation_report.json
\end{DocCode}

A compatibility exporter may optionally reconstruct skipped children for old validators that expect them.
If skipped children are emitted, they must follow the historical null convention:

\begin{DocCode}
{
  "parameters": {
    "<parameter>": null
  },
  "parameter_hash": null,
  "time0_baseline_uuid": null,
  "time0_baseline_hash": null,
  "direction": "increase",
  "intervention_time": null
}
\end{DocCode}

Skipped children must never be materialized by Stage: Simulate.
They are coverage annotations only.

\section*{Validation Rules for Compatibility Exports}

A generated compatibility file is valid only if:

\begin{itemize}
\item every runnable baseline entry maps to one baseline node in \code{run_nodes.jsonl};
\item every runnable child entry maps to one intervention node in \code{run_nodes.jsonl};
\item every non-null \code{time0_baseline_uuid} maps to one time-zero/static-change node in \code{run_nodes.jsonl};
\item every runnable child changes exactly one exposed root-cause parameter;
\item every runnable child has an \code{intervention_time} satisfying \code{0 < intervention_time < experiment_config.json::end_time_input_s};
\item every emitted identifier is unique within the compatibility file;
\item every skipped child, if present, uses the historical null convention;
\item no compatibility entry changes the canonical recipe represented by the resolved run graph.
\end{itemize}

\section*{Artifact Policy}

\begin{itemize}
\item \code{model_run_specs.json} is no longer authored or modified by humans.
\item The resolved run graph is the canonical planning artifact.
\item \code{model_run_specs.json} may be generated only as a legacy compatibility artifact.
\item Manual edits to \code{model_run_specs.json} are discouraged because they can diverge from \code{run_nodes.jsonl} and \code{run_edges.jsonl}.
\item If the compatibility file is stale, regenerate it from the resolved run graph.
\item Stage: Simulate should consume \code{run_nodes.jsonl} and \code{run_edges.jsonl}; it should not refresh or repair \code{model_run_specs.json}.
\end{itemize}

\section*{Migration Guidance}

\begin{center}
\begin{tabular}{p{0.34\linewidth}p{0.56\linewidth}}
\hline
\textbf{Old dependency} & \textbf{New dependency} \\
\hline
\code{baseline_uuid} & baseline \code{run_id} from \code{run_nodes.jsonl} \\
\code{intervention_uuid} & intervention \code{run_id} from \code{run_nodes.jsonl} \\
\code{time0_baseline_uuid} & time-zero/static-change \code{run_id} from \code{run_nodes.jsonl} \\
\code{baseline_parameters_hash} & baseline node \code{recipe_hash} \\
\code{parameter_hash} & intervention node \code{recipe_hash} \\
\code{time0_baseline_hash} & time-zero/static-change node \code{recipe_hash} \\
nested \code{children} object & \code{run_edges.jsonl} graph relationships \\
skipped null children & \code{generation_report.json} coverage report \\
\hline
\end{tabular}
\end{center}

Once all downstream stages consume the resolved run graph directly, this compatibility document and the compatibility exporter can be removed.

\end{document}
